\documentclass[11pt,a4paper]{article}

\usepackage[T1]{fontenc}
\usepackage[utf8]{inputenc}
\usepackage[ngerman]{babel}
\usepackage{amsmath,amssymb}
\usepackage{graphicx}
\usepackage{booktabs}
\usepackage{geometry}
\usepackage{xcolor}
\usepackage{listings}
\usepackage{pgfplots}
\pgfplotsset{compat=1.18}

\geometry{a4paper, margin=2.5cm}

\definecolor{codebg}{rgb}{0.95,0.95,0.95}
\lstset{
    backgroundcolor=\color{codebg},
    basicstyle=\ttfamily\small,
    keywordstyle=\color{blue},
    commentstyle=\color{gray},
    stringstyle=\color{red!60!black},
    numbers=left,
    numberstyle=\tiny\color{gray},
    frame=single,
    breaklines=true,
    showstringspaces=false,
    language=Python,
    literate=%
        {ä}{{\"a}}1 {ö}{{\"o}}1 {ü}{{\"u}}1
        {Ä}{{\"A}}1 {Ö}{{\"O}}1 {Ü}{{\"U}}1
        {ß}{{\ss}}1
}

\title{Numerische Lösung der eindimensionalen Wärmeleitungsgleichung\\
mittels expliziter finiter Differenzen}
\author{Max\ Musterman \\ \small Institut für Angewandte Mathematik, Musterstadt}
\date{\today}

\begin{document}
\maketitle

\begin{abstract}
In dieser Arbeit wird die eindimensionale Wärmeleitungsgleichung mit einem expliziten
Finite-Differenzen-Verfahren numerisch gelöst. Wir leiten das Diskretisierungsschema her,
diskutieren die Stabilitätsbedingung nach von Neumann und implementieren das Verfahren
in Python. Die numerischen Ergebnisse werden mit der analytischen Lösung verglichen und
die Konvergenzordnung des Verfahrens wird empirisch bestätigt.
\end{abstract}

\section{Einleitung}
Die Wärmeleitungsgleichung ist eine der grundlegenden partiellen Differentialgleichungen
der mathematischen Physik. Sie beschreibt die zeitliche Entwicklung der Temperaturverteilung
$u(x,t)$ in einem Medium und lautet in ihrer eindimensionalen Form
\begin{equation}
\frac{\partial u}{\partial t} = \alpha \, \frac{\partial^2 u}{\partial x^2}, \qquad x \in [0, L], \; t > 0,
\label{eq:heat}
\end{equation}
wobei $\alpha > 0$ die Temperaturleitfähigkeit des Materials bezeichnet. Zusammen mit den
Randbedingungen
\begin{equation}
u(0,t) = u(L,t) = 0
\end{equation}
und der Anfangsbedingung $u(x,0) = u_0(x)$ ist das Problem wohlgestellt. Obwohl für viele
Anfangsbedingungen eine analytische Lösung mittels Fourier-Reihen existiert, ist für
komplexere Geometrien oder variable Koeffizienten $\alpha(x)$ ein numerisches Verfahren
notwendig. Im Folgenden betrachten wir das einfachste derartige Verfahren: die explizite
Finite-Differenzen-Methode (FTCS, \emph{forward-time centered-space}).

\section{Diskretisierung}
Wir unterteilen das Ortsintervall $[0, L]$ in $N$ gleich große Teilintervalle der Länge
$\Delta x = L / N$ und die Zeitachse in Schritte der Länge $\Delta t$. Die Gitterpunkte
seien $x_i = i \, \Delta x$ für $i = 0, \dots, N$ und $t_n = n \, \Delta t$. Wir approximieren
$u(x_i, t_n) \approx u_i^n$.

Die Zeitableitung wird durch einen Vorwärtsdifferenzenquotienten ersetzt,
\begin{equation}
\frac{\partial u}{\partial t}\bigg|_{i}^{n} \approx \frac{u_i^{n+1} - u_i^n}{\Delta t},
\end{equation}
und die räumliche zweite Ableitung durch den zentralen Differenzenquotienten
\begin{equation}
\frac{\partial^2 u}{\partial x^2}\bigg|_{i}^{n} \approx \frac{u_{i+1}^n - 2 u_i^n + u_{i-1}^n}{\Delta x^2}.
\end{equation}
Einsetzen in Gleichung~\eqref{eq:heat} liefert das explizite Update-Schema
\begin{equation}
u_i^{n+1} = u_i^n + r \left( u_{i+1}^n - 2 u_i^n + u_{i-1}^n \right), \qquad
r := \frac{\alpha \, \Delta t}{\Delta x^2}.
\label{eq:update}
\end{equation}
Eine von-Neumann-Stabilitätsanalyse zeigt, dass das Schema~\eqref{eq:update} genau dann
stabil ist, wenn
\begin{equation}
r = \frac{\alpha \, \Delta t}{\Delta x^2} \leq \frac{1}{2}
\label{eq:stability}
\end{equation}
gilt. Diese Bedingung schränkt den Zeitschritt $\Delta t$ bei feiner Ortsauflösung
erheblich ein, da $\Delta t \propto \Delta x^2$ gefordert wird.

\section{Implementierung}
Das Verfahren aus Gleichung~\eqref{eq:update} lässt sich direkt in Python implementieren.
Der folgende Code zeigt eine vektorisierte Umsetzung mit \texttt{NumPy}, welche die
Randbedingungen bei jedem Zeitschritt erzwingt:

\begin{lstlisting}[caption={Explizites FTCS-Verfahren für die Wärmeleitungsgleichung}]
import numpy as np

def loese_waermeleitung(u0, alpha, dx, dt, n_schritte):
    """Loest die 1D-Waermeleitungsgleichung explizit."""
    r = alpha * dt / dx**2
    if r > 0.5:
        raise ValueError(f"Instabil: r = {r:.3f} > 0.5")

    u = u0.copy()
    for n in range(n_schritte):
        u_neu = u.copy()
        u_neu[1:-1] = u[1:-1] + r * (u[2:] - 2*u[1:-1] + u[:-2])
        u_neu[0] = 0.0   # Dirichlet-Randbedingung
        u_neu[-1] = 0.0
        u = u_neu
    return u
\end{lstlisting}

Die Stabilitätsbedingung wird in Zeile 6 explizit geprüft, sodass instabile
Parameterkombinationen frühzeitig abgefangen werden.

\section{Ergebnisse}
Als Testfall wählen wir $L = 1$, $\alpha = 1$ und die Anfangsbedingung
$u_0(x) = \sin(\pi x)$, für die die analytische Lösung
\begin{equation}
u_{\text{exakt}}(x,t) = e^{-\pi^2 \alpha t} \sin(\pi x)
\end{equation}
bekannt ist. Abbildung~\ref{fig:profile} zeigt die numerische Lösung zu verschiedenen
Zeitpunkten im Vergleich zur analytischen Lösung.

\begin{figure}[htbp]
\centering
\begin{tikzpicture}
\begin{axis}[
    width=0.8\textwidth,
    height=6cm,
    xlabel={$x$},
    ylabel={$u(x,t)$},
    legend pos=north east,
    grid=major,
]
\addplot[domain=0:1, samples=100, blue, thick] {sin(deg(pi*x))*exp(-pi^2*0.0)};
\addlegendentry{$t = 0.00$}
\addplot[domain=0:1, samples=100, red, thick, dashed] {sin(deg(pi*x))*exp(-pi^2*0.05)};
\addlegendentry{$t = 0.05$}
\addplot[domain=0:1, samples=100, green!60!black, thick, dotted] {sin(deg(pi*x))*exp(-pi^2*0.15)};
\addlegendentry{$t = 0.15$}
\end{axis}
\end{tikzpicture}
\caption{Zeitliche Entwicklung des Temperaturprofils $u(x,t)$ für $\alpha = 1$.
Die Amplitude klingt exponentiell mit der Rate $\pi^2 \alpha$ ab.}
\label{fig:profile}
\end{figure}

Um die Konvergenzordnung des Verfahrens zu überprüfen, berechnen wir den maximalen
Fehler $E_N = \max_i |u_i^{n_{\text{end}}} - u_{\text{exakt}}(x_i, t_{\text{end}})|$
für verschiedene Ortsauflösungen $N$ bei fester Stabilitätszahl $r = 0.4$.
Tabelle~\ref{tab:convergence} fasst die Ergebnisse zusammen.

\begin{table}[htbp]
\centering
\caption{Maximaler Fehler und experimentelle Konvergenzordnung $p$ in Abhängigkeit von $N$.}
\label{tab:convergence}
\begin{tabular}{@{}rccc@{}}
\toprule
$N$ & $\Delta x$ & $E_N$ & $p = \log_2(E_{N/2}/E_N)$ \\
\midrule
20  & $5.00 \times 10^{-2}$ & $8.21 \times 10^{-3}$ & --   \\
40  & $2.50 \times 10^{-2}$ & $2.05 \times 10^{-3}$ & 2.00 \\
80  & $1.25 \times 10^{-2}$ & $5.13 \times 10^{-4}$ & 2.00 \\
160 & $6.25 \times 10^{-3}$ & $1.28 \times 10^{-4}$ & 2.00 \\
\bottomrule
\end{tabular}
\end{table}

Wie erwartet zeigt das Verfahren die theoretisch vorhergesagte Konvergenzordnung
$p \approx 2$ in $\Delta x$, da sowohl der zentrale Differenzenquotient in Gleichung~(4)
als auch die Zeitintegration bei konstantem $r$ quadratische Konvergenz liefern.

\section{Zusammenfassung}
Wir haben das explizite FTCS-Verfahren zur Lösung der Wärmeleitungsgleichung
hergeleitet, implementiert und numerisch validiert. Die von-Neumann-Stabilitätsbedingung
$r \leq 1/2$ erwies sich als notwendig und hinreichend für stabile Lösungen, und die
empirische Konvergenzordnung stimmt mit der theoretischen Erwartung von $\mathcal{O}(\Delta x^2)$
überein. Für praktische Anwendungen mit feiner Ortsauflösung ist jedoch aufgrund der
Bedingung $\Delta t \propto \Delta x^2$ ein implizites Verfahren, etwa das
Crank-Nicolson-Schema, dem hier vorgestellten expliziten Ansatz vorzuziehen, da es
unbedingt stabil ist und größere Zeitschritte erlaubt.

\end{document}
